\documentclass{standalone}
\usepackage{graphicx}	
\usepackage{amssymb, amsmath}
\usepackage{color}

\usepackage{tikz}
\usetikzlibrary{intersections, backgrounds, math, arrows.meta}
\usepackage{pgfmath}

\definecolor{light}{RGB}{220, 188, 188}
\definecolor{mid}{RGB}{185, 124, 124}
\definecolor{dark}{RGB}{143, 39, 39}
\definecolor{highlight}{RGB}{180, 31, 180}
\definecolor{gray10}{gray}{0.1}
\definecolor{gray20}{gray}{0.2}
\definecolor{gray30}{gray}{0.3}
\definecolor{gray40}{gray}{0.4}
\definecolor{gray60}{gray}{0.6}
\definecolor{gray70}{gray}{0.7}
\definecolor{gray80}{gray}{0.8}
\definecolor{gray90}{gray}{0.9}
\definecolor{gray95}{gray}{0.95}

\tikzmath{
  function marginal(\x) {
    \a=3;
    \b=3;
    return exp( (\a - 1) * ln(\x) + (\b - 1) * ln(1 - \x) + 3.401197);
  };
}

\begin{document}

\begin{tikzpicture}[scale=1]

  \begin{scope}[shift={(0, 0)}]
    \draw[white] (-4.5, -3) rectangle (3.5, 2.5);
      
    \begin{scope}
      \clip (-3, -2.1) rectangle (3, 2);
      \draw[dark, line width=1] (-3, -2) --
        plot[domain=0.001:0.999, smooth, samples=100, variable=\x] 
          ( { 6 * \x - 3 }, { 1.75 * marginal(\x) - 2 } )
        -- (3, 2);
    \end{scope}
    
    \draw [->, >=stealth, line width=1] (-3.00, -2.015) -- +(0, 4);
    \draw [->, >=stealth, line width=1] (+3.00, -2.015) -- +(0, 4);
    \draw [line width=1] (-3.015, -2.00) -- +(6, 0);
    
    \draw[line width=1] (-3, -2) -- (-3, -2.1);
    \node at (-3, -2.25) { $0$ };

    \draw[line width=1] (+3, -2) -- (+3, -2.1);
    \node at (+3, -2.25) { $1$ };
    
    \node at (-3.75, 0) { $p( x_{1} )$ };
    \node at (0, -2.75) { $x_{1}$ };
    
  \end{scope}

\end{tikzpicture}

\end{document}  